\documentclass[11pt]{beamer}
\usetheme{metropolis}
\usepackage{appendixnumberbeamer}
\usepackage[T1]{fontenc}
\usepackage{listings}
\usepackage{xcolor}

\setbeamertemplate{frame numbering}[fraction]

\lstset{
  basicstyle=\ttfamily\footnotesize,
  keywordstyle=\color{mDarkTeal}\bfseries,
  commentstyle=\color{gray},
  stringstyle=\color{mDarkBrown},
  showstringspaces=false,
  breaklines=true,
  columns=fullflexible,
  language=Python,
}

\title{Cookiecutter Data Science Conventions}
\subtitle{How this repo keeps agents honoring project structure}
\author{Josh Goldman}
\date{\today}

\begin{document}

\maketitle

\begin{frame}{What we built}
  Three pieces, working together:
  \begin{enumerate}
    \item \textbf{\texttt{CLAUDE-CCDS.md}} --- a conventions file that tells
      an agent \emph{how to write code} in a ccds-structured repo, day to
      day (not just at scaffold time): directory layout, config/CLI
      conventions, \texttt{Protocol}-based design, testing, tooling.
    \item \textbf{\texttt{new-ds-project} skill} --- scaffolds a
      \textbf{brand-new} repo into this structure from scratch.
    \item \textbf{\texttt{refactor-ccds} skill} --- takes an
      \textbf{existing}, non-ccds repo and reorganizes it into our style,
      rather than moving files ad hoc.
  \end{enumerate}
  \vspace{0.5em}
  \small The two skills handle one-time setup/migration; the \texttt{.md}
  file keeps every later change consistent with it.
\end{frame}

\begin{frame}[fragile]{ccds file structure}
  \begin{lstlisting}[language=,basicstyle=\ttfamily\scriptsize]
data/
    raw/          # immutable, gitignored, never edited in place
    interim/      # regenerable from raw
    processed/    # regenerable, model-ready
notebooks/        # exploration only: <n>.<initials>-<desc>.ipynb
models/           # serialized trained models
reports/
    figures/      # generated charts
references/       # data dictionaries, explanatory material
docs/adr/         # architecture decision records
src/<module_name>/
    config.py         # paths, env, logging, seeding
    data/               make_dataset.py, validators.py
    features/*          base.py, build_features.py
    modeling/            base.py, train.py, predict.py, architectures.py
    visualization/        plots.py
tests/            # mirrors src/<module_name>/, one test file per source file
\end{lstlisting}
\end{frame}

\begin{frame}{Tools we're using}
  \begin{columns}[T]
    \column{0.5\textwidth}
    \textbf{Config \& CLI}
    \begin{itemize}
      \item \texttt{pathlib} + \texttt{config.py}
      \item \texttt{python-dotenv}
      \item \texttt{typer} (simple stages)
      \item \texttt{Tyro} (config-heavy stages)
      \item \texttt{loguru} + \texttt{tqdm}
    \end{itemize}
    \vspace{0.5em}
    \textbf{Data \& experiments}
    \begin{itemize}
      \item \texttt{pandera} (schema validation)
      \item \texttt{wandb} (run tracking)
      \item \texttt{Makefile} (orchestration)
    \end{itemize}
    \column{0.5\textwidth}
    \textbf{Quality}
    \begin{itemize}
      \item \texttt{ruff} (lint + format)
      \item \texttt{mypy} (typing)
      \item \texttt{pytest} (tests)
      \item \texttt{pre-commit} (blocks commits)
      \item \texttt{pip-audit} (dependency security)
    \end{itemize}
  \end{columns}
\end{frame}

\begin{frame}[fragile]{Method structure as \texttt{Protocol}s}
  \small
  \begin{itemize}
    \item Each stage's method structure is broken into a
      \texttt{typing.Protocol} interface with swappable concrete
      implementations --- decomposed into the \emph{real} sub-parts rather
      than one interface per stage:
  \end{itemize}
\begin{lstlisting}[basicstyle=\ttfamily\scriptsize]
class FeatureExtractor(Protocol): def transform(self, raw): ...
class Backbone(Protocol):         def encode(self, features): ...
class PredictorHead(Protocol):    def forward(self, embedding): ...
class Model(Protocol):
    def fit(self, features, labels) -> None: ...
    def predict(self, features): ...
\end{lstlisting}
  \begin{itemize}
    \item This decomposition also gives us something to
      \textbf{visualize}: diagram the protocols and their
      information flow/containment (what calls what, what wraps what) to
      see a program's high-level structure and building blocks at a
      glance --- independent of any one implementation's internals.
  \end{itemize}
\end{frame}

\begin{frame}{Style \& workflow --- the small things}
  \begin{itemize}
    \item \textbf{Style}: \texttt{ruff} enforces PEP-8-ish formatting +
      linting, roughly followed rather than dogmatically; NumPy-style
      docstrings on public functions/classes; comments explain \emph{why},
      not \emph{what}; no commented-out code or context-free TODOs.
    \item \textbf{Git/GitHub}:
      \begin{itemize}
        \small
        \item Commit early and often, imperative-mood messages explaining
          \emph{why}.
        \item Branch for anything experimental/risky; merge or discard.
        \item Push regularly --- it's the backup and shareable record.
        \item Never force-push or rewrite history on \texttt{main}.
        \item No PR requirement or branch-protection rules by default.
      \end{itemize}
  \end{itemize}
\end{frame}

\begin{frame}{How \texttt{refactor-ccds} works}
  \small
  Never move files ad hoc --- analyze before touching anything:
  \begin{enumerate}
    \item \textbf{Inventory \& analyze} the repo first: classify every
      file, flag anything that \textbf{doesn't cleanly fit} the standard
      structure or would need an actual behavior change to move.
    \item \textbf{Grill session} on that ambiguous set --- resolve every
      open question with the user before moving a single file.
    \item \textbf{Prepare a clear plan}: full old $\rightarrow$ new path
      mapping, deletions, and behavior changes called out explicitly.
    \item \textbf{Present for approval} --- wait for explicit confirmation
      before touching the filesystem.
    \item \textbf{Then refactor}: move on a new branch with \texttt{git mv}
      (history preserved), fix what broke, check tests against baseline.
  \end{enumerate}
\end{frame}

\begin{frame}{Future steps}
  \begin{itemize}
    \item Right now, \texttt{CLAUDE-CCDS.md} bakes in \textbf{one fixed}
      set of style choices (Tyro vs. typer, \texttt{Protocol} decomposition,
      wandb, \dots) --- the same for every project and every user.
    \item \textbf{Make it customizable}: turn these choices into
      preferences instead of hardcoded defaults, so different
      projects/collaborators can deviate where it makes sense.
    \item Concretely: a \textbf{skill that interviews the user} about their
      preferred style (naming, tooling choices, how strict to be about
      \texttt{Protocol}s, testing rigor, \dots), then
      \textbf{regenerates/updates} the project's \texttt{CLAUDE.md}
      conventions from those answers --- rather than everyone hand-editing
      the same fixed file.
  \end{itemize}
\end{frame}

\end{document}
